% !TEX root = ../main.tex
\chapter{CWT、RCWA、FDTD、FEM 的比较与互证}
\label{ch:method-compare}

\section*{学习目标}
\begin{itemize}
  \item 清楚区分 CWT、RCWA、FDTD、FEM 各自解决的物理问题，而不是只记住它们的软件标签。
  \item 建立多方法互证的意识，理解为什么单一方法几乎不可能覆盖 PCSEL 的全部核心问题。
  \item 掌握一条从概念设计到器件级验证的推荐工作流，并知道每一层输出该如何交给下一层。
\end{itemize}

\section*{先修知识}
建议已掌握 \cref{ch:cwt-min,ch:cwt-3d,ch:rcwa,ch:fdtd-fem} 中关于半解析模型、开放周期结构和全波有限器件求解的内容。

\section*{本章结构}
\ChapterModelLevel{L1→L3}
到这里为止，书里已经出现了四类常用工具：CWT、RCWA、FDTD 和 FEM。初学者常见错误，不是不会用它们，而是误以为其中一种就足够覆盖全部问题。本章的任务就是把这四种方法放到同一张地图上，先比较它们分别擅长什么，再说明它们的互证原则，最后给出一条真正适合 PCSEL 的研究工作流。这样安排不是为了增加方法数量，而是为了避免结论层级混乱。

\section{为什么单一方法几乎注定不够}
PCSEL 同时是三种对象：它是周期开放电磁系统，是半导体有源腔，也是带真实注入与热分布的器件。任何一种方法，只要天然擅长其中一层，就通常会在另外两层变弱。

CWT 擅长把复杂物理压缩成可解释的耦合通道，因此适合建立直觉、做灵敏度分析和前期筛选；但它的代价是近似较强，不能天然承担高保真几何与开放边界的全部负担。

RCWA 擅长周期开放单胞，特别适合研究辐射通道、偏振和开放共振；但它默认的是无限周期单胞，不天然处理有限器件边缘。

FDTD 和 FEM 擅长真实有限器件全波分析，能把边缘、复杂层栈和真实几何带回来；但它们计算代价更高，也通常不像 CWT 那样自动给出清晰的机制分解。

于是，单一方法不够不是因为“每种方法都不行”，而是因为 PCSEL 的问题本来就不是单层问题。方法必须分层合作，才和物理对象的层次相匹配。

\begin{IntuitionBox}
成熟的 PCSEL 建模不应寻找“最强求解器”，而应建立一条从粗到细、从机理到验证、从光学到器件的层级链。方法之间不是竞赛关系，而是接力关系。
\end{IntuitionBox}

\section{先按“问题类型”而不是“方法名字”来比较}
最有用的比较方式，不是问“哪个方法更准”，而是先问“我现在究竟在回答什么问题”。如果问题是候选模目录、主导耦合通道和几何灵敏度，那么 CWT 最合适；如果问题是开放周期单胞的辐射、反射谱和光学阈值增益趋势，RCWA 最合适；如果问题是有限器件边缘、近远场、真实 $Q$ 和复杂几何，FDTD/FEM 更合适；如果问题进一步变成阈值电流、热稳定性和模式切换，则必须把这些光学结果交给器件模型。

也就是说，方法比较的第一原则是“按问题选择方法”，而不是“按方法反推问题”。这条原则一旦建立，许多常见争论其实可被重新表述。比如，“RCWA 准还是 FDTD 准”这种问题本身就不完整，因为两者默认解决的问题层级并不一样。

\begin{table}[htbp]
  \centering
  \small
  \caption{几类常用方法的典型计算量级与适用尺度。表中时间与自由度属于\textbf{[本书算例]}，仅用于排程直觉，不用于跨软件绝对比较。}
  \label{tab:method_cost_scale}
  \begin{tabularx}{\textwidth}{@{}p{1.8cm}p{1.9cm}p{2.4cm}Y@{}}
    \toprule
    方法 & 典型自由度 & 典型对象尺度 & 单次求解量级 \\ 
    \midrule
    PWEM & $\sim10^2$--$10^3$ & 二维单胞、被动无限周期 & 秒级到分钟级 \\
    CWT & $\sim10^0$--$10^1$ & 少数主导波与参数扫描 & 秒级 \\
    RCWA & $\sim10^3$--$10^4$ & 三维层状周期单胞 & 秒级到分钟级 \\
    3D FDTD & $\sim10^7$--$10^8$ & 数十周期有限孔阵、宽频瞬态 & 小时级到天级 \\
    3D FEM & $\sim10^6$--$10^7$ & 有限器件、本征模或窄频响应 & 小时级到天级 \\
    \bottomrule
  \end{tabularx}
\end{table}

\textbf{benchmark 前提（用于解释表~\ref{tab:method_cost_scale}）}：
以单机 16 核 CPU、64 GB 内存、双精度求解为参考，RCWA 与 FDTD/FEM 的网格/谐波截断分别取“达到章节内误差门限”的最小设置；若改用 GPU、不同求解器实现或更严格收敛标准，耗时可变化一个数量级以上。因此表中“秒级/小时级/天级”应读作\textbf{[工程经验]}区间，而非方法学常数。RCWA 资源量级对截断阶数与因子化规则的敏感性，可参见经典收敛文献\cite{moharam1995stablercwa,li1996}。

把 \cref{tab:method_cost_scale} 放在这里的目的，不是制造“谁更高级”的错觉，而是说明工作流必须分层：不同方法面对的物理问题不同，计算代价也相差数个数量级。若把本来应由 PWEM 或 RCWA 完成的候选筛选全部交给 3D FDTD/FEM，常见后果不是更严谨，而是更慢，也更难定位错误。
为了帮助读者更直观地理解方法分工，\cref{fig:ch15_method_comparison_radar}给出一个六维雷达\textbf{[教学示意]}图。图中多边形仅表达“优势方向”和“取舍关系”，不代表跨软件、跨硬件可复现实测分数；若要做发布级定量比较，必须同时给出 benchmark 四元信息（硬件、求解器版本、几何与边界、误差门槛）以及原始数据表。

% !TEX root = ../main.tex
% Figure: problem-first method map for PCSEL simulation

\begin{figure}[htbp]
  \centering
  \resizebox{\textwidth}{!}{%
  \begin{tikzpicture}[x=1cm,y=1cm,>=Latex,font=\small]
    \tikzset{
      question/.style={draw=black!60,rounded corners=2pt,fill=gray!6,align=center,minimum width=3.2cm,minimum height=0.78cm,font=\footnotesize},
      method/.style={draw=blue!55!black,rounded corners=2pt,fill=blue!7,align=center,minimum width=2.75cm,minimum height=0.82cm,font=\bfseries\footnotesize},
      output/.style={draw=green!45!black,rounded corners=2pt,fill=green!7,align=center,text width=3.15cm,minimum height=1.05cm,font=\footnotesize},
      caution/.style={draw=red!60!black,rounded corners=2pt,fill=red!6,align=left,text width=4.45cm,font=\footnotesize},
      flow/.style={->,thick},
      softflow/.style={->,thick,dashed,gray!75}
    }

    \node[font=\bfseries] at (7.5,5.25) {先按问题类型选方法，而不是按软件名字倒推结论};

    \node[question] (q1) at (1.8,4.1) {候选频段在哪里?};
    \node[method] (m1) at (1.8,2.95) {PWEM / 能带};
    \node[output] (o1) at (1.8,1.55) {频率窗口\\简并关系\\高对称点标签};

    \node[question] (q2) at (5.6,4.1) {主导耦合是谁?};
    \node[method] (m2) at (5.6,2.95) {CWT / CMT};
    \node[output] (o2) at (5.6,1.55) {散射通道\\参数灵敏度\\模式分裂方向};

    \node[question] (q3) at (9.4,4.1) {开放单胞如何辐射?};
    \node[method] (m3) at (9.4,2.95) {RCWA};
    \node[output] (o3) at (9.4,1.55) {上下出光\\辐射损耗\\偏振趋势};

    \node[question] (q4) at (13.2,4.1) {有限器件是否仍成立?};
    \node[method] (m4) at (13.2,2.95) {FDTD / FEM};
    \node[output] (o4) at (13.2,1.55) {有限阵列近远场\\边缘损耗\\真实层栈修正};

    \foreach \a/\b in {q1/m1,m1/o1,q2/m2,m2/o2,q3/m3,m3/o3,q4/m4,m4/o4} {
      \draw[flow] (\a) -- (\b);
    }
    \foreach \a/\b in {o1/q2,o2/q3,o3/q4} {
      \draw[softflow] (\a.east) -- ++(0.35,0) |- (\b.west);
    }

    \node[caution,anchor=north west] at (0.25,0.25) {
      \textbf{层级边界：}\\
      PWEM/CWT 不能单独给阈值电流；\\
      RCWA 不能单独给有限阵列边缘排序；\\
      冷腔 FDTD/FEM 不能单独给热-电-光工作窗口。
    };
    \node[draw=black!45,rounded corners=2pt,fill=yellow!10,align=center,text width=8.6cm,font=\footnotesize,anchor=north west] at (5.35,0.25) {
      最可靠的结论通常来自“同一物理解释链”在多个层级上保持一致：\\
      候选模身份稳定、对称性标签稳定、参数趋势稳定、损耗来源可追踪。
    };
  \end{tikzpicture}%
  }
  \caption{PCSEL 数值方法的分工地图。图中强调从问题类型进入：能带/PWEM 给候选目录，CWT/CMT 给机理压缩，RCWA 给开放周期单胞辐射，FDTD/FEM 给有限器件与真实层栈验证。它是教学示意图，不代表跨软件 benchmark 分数。}
  \label{fig:ch15_method_comparison_radar}
\end{figure}

\section{四种方法各自最擅长什么}
为了把分工说得更清楚，可以用一句话概括四种方法的适用场景。

CWT 最擅长回答：“为什么这类几何扰动会主要改这个耦合通道，而不是另一个？”

RCWA 最擅长回答：“这个周期单胞在开放边界下会怎样把能量耦合到各个辐射通道？”

FDTD 最擅长回答：“这个有限器件在宽频和瞬态条件下，实际会怎样建立、衰减和辐射？”

FEM 最擅长回答：“这个复杂几何、复杂层栈的有限器件在某个频率或某支模式附近，精确场分布和复本征值是什么？”

如果把四者的适用场景混在一起用，就会出现常见误判。例如：用 CWT 试图给出精确远场角宽，用 RCWA 试图给出有限器件边缘模泄漏，用 FDTD 的一组冷腔结果直接断言器件阈值电流。这些都不是方法本身的问题，而是问题与方法没有对齐。

\begin{table}[htbp]
  \centering
  \footnotesize
  \caption{各方法最常见的“明确不适用”场景。表中“不适用”指不应让该方法单独承担最终结论。}
  \label{tab:method_not_applicable}
  \begin{tabularx}{\textwidth}{@{}p{1.9cm}Y Y@{}}
    \toprule
    方法 & 不应单独承担的结论 & 原因 \\ \midrule
    CWT & 精确远场角宽、有限阵列边缘泄漏、器件级阈值电流 & 基底维数过低，且有限尺寸与电热回写未显式进入 \\
    PWEM & 开放辐射损耗、真实有限器件排序、阈值结论 & 典型形式是无损被动无限周期本征问题 \\
    RCWA & 有限孔阵边缘主导问题、横向包络与边缘吸收主导问题 & 单胞周期模型本身不含真实有限尺寸终止 \\
    FDTD/FEM 冷腔 & 电注入阈值电流、阈上稳态、热稳定工作窗口 & 若无电热回写，只能给冷腔或小信号光学层级结论 \\
    \bottomrule
  \end{tabularx}
\end{table}

\section{比较方法时，应比较哪些量}
多方法互证并不意味着所有数值都要逐位相同。不同方法的基底、边界和离散方式不同，过分追求“完全一致”往往会掩盖更重要的信息。更有意义的是比较那些应当在方法变化下保持稳定的物理量或趋势。

第一类是候选模的排序。例如目标频段内究竟是哪几支模式最值得关注，这种排序不应因方法切换而完全翻转。

第二类是对称性与偏振标签。虽然数值细节可能不同，但目标模属于哪类对称性表示、主要偏向哪种偏振，应该能在方法间对得上。

第三类是参数微扰趋势。若增大孔半径在 CWT 中预测会增强某种辐射通道，那么 RCWA 与全波方法至少应在趋势上给出相容结论。

第四类是损耗来源分解。即使绝对 $Q$ 不完全一致，也应能判断损耗主要来自向上辐射、向下辐射、边缘泄漏还是掺杂吸收重叠。

换句话说，互证的重点不是“每个小数点都一样”，而是“物理解释链条在不同分辨率下是否一致”。

\begin{RigorousBox}
跨方法比较时，最危险的不是少量数值偏差，而是层级混淆。若一种方法给的是无限周期被动结果，而另一种方法给的是有限器件冷腔结果，再拿二者的绝对阈值优劣直接比较，本身就没有严格可比性。先对齐问题层级，再比较结果，才有意义。
\end{RigorousBox}

\section{数值误差从哪来：互证前的最后一项功课}
每种方法都有自己的主导误差源与收敛行为：PWEM 的平面波截断与有效折射率近似、CWT 的模型降阶、RCWA 的 Fourier 收敛与 Gibbs 现象、FDTD 的时域误差累积、FEM 的网格与弱形式实现。逐方法的误差源清单、三层收敛门槛口径（探索级/本书算例/发布级）、不确定度传递与报告规范，统一收入 \cref{ch:sim-appendix} 的数值不确定度部分（\cref{sec:numerical-error-per-method}）。互证时只需带着一条纪律：任何进入比较的数字，都必须附有收敛证据和误差量级。

\section{为什么互证必须按“由粗到细”的顺序组织}
如果一开始就直接上三维全器件全波扫描，当然不是做不到，但通常代价极高，而且一旦结果不理想，很难知道问题究竟出在晶格、外延还是边界设置。相反，若先由粗到细组织工作流，每一层都能承担一个清楚的筛选任务，后续搜索空间会大幅缩小。

最自然的顺序通常是：先用 \cref{ch:gamma} 的能带阅读框架配合 \cref{ch:pwe} 的 PWEM 找到候选模，再用 CWT 识别主导耦合通道与参数方向，再用 RCWA 精修单胞开放散射与光学阈值趋势，再用 FDTD/FEM 把有限尺寸和真实外延层栈带回来，最后把这些结果交给电注入和热耦合模型。这条顺序并不是人为规定，而是与问题层级天然匹配。

这样做的另一个好处是，任何阶段一旦出问题，都能更容易定位。例如，如果 CWT 和 RCWA 已经都显示某方案目标模并无明显优势，那就没必要花大代价做全器件验证；反过来，如果 RCWA 很乐观、但全波有限器件结果很差，问题通常就出在边缘、有限尺寸或纵向层栈细节，而不是单胞本身。

\section{一条适合 PCSEL 的推荐工作流}
把前面的思路压缩成一条实际可执行的工作流，可以概括为五步。

第一步，建立候选模地图。使用 \cref{ch:gamma} 的能带概念和 \cref{ch:pwe} 的 PWEM 计算，在目标波长附近找出 Gamma 点模族、简并关系和对称性标签。

第二步，做机理级预设计。使用 CWT 或更高阶 CWT，把注意力集中到少数关键耦合通道，判断哪些几何扰动最可能改变模式分裂、偏振选择和辐射窗口。

第三步，做单胞开放精修。使用 RCWA，围绕前两步锁定的候选模区间，提取被动谱、上下辐射分配和光学阈值增益趋势。

第四步，做有限器件全波验证。使用 FDTD/FEM，把有限孔阵边缘、真实层栈、窗口和边界条件带进模型，检查真实近场、远场、有限尺寸 $Q$ 和模式排序。

第五步，做器件级回写链。把第四步输出的模式重叠、损耗和导数信息交给增益模型、电学输运与热耦合模型，得到阈值电流、工作点漂移和高功率稳定性判断。

这条工作流最大的价值，不在于“把所有工具都用一遍”，而在于每一步都有清楚的淘汰与交接标准。

\begin{table}[htbp]
  \centering
  \footnotesize
  \caption{方法链中的交接表：每一层最值得交给下一层的中间量。}
  \label{tab:method_handoff}
  \begin{tabularx}{\textwidth}{@{}p{2.1cm}p{3.1cm}Y@{}}
    \toprule
    当前层 & 主要中间量 & 下一层如何使用 \\ \midrule
    PWEM & 候选频段、Gamma 点模目录、简并与对称性标签 & 交给 CWT 选择最小耦合子空间，交给 RCWA 锁定扫描频段 \\
    CWT & 主导耦合通道、模式分裂方向、参数灵敏度符号 & 交给 RCWA 选重点参数和候选模，减少盲扫 \\
    RCWA & 单胞开放谱、上下辐射分配、光学阈值代理、偏振趋势 & 交给 FDTD/FEM 构造有限阵列验证目标，交给器件模型提供辐射损耗输入 \\
    FDTD/FEM & 有限器件近场/远场、有限尺寸 $Q$、场重叠与边缘泄漏 & 交给增益、电学和热模型作为器件级输入 \\
    电热光模型 & 阈值电流、温升、工作点漂移、稳定工作窗口 & 交给实验回写链做回验与设计更新 \\
    \bottomrule
  \end{tabularx}
\end{table}

\begin{table}[htbp]
  \centering
  \footnotesize
  \caption{工作流的继续/停机判据。若未通过当前层门槛，不应升级到下一层最昂贵模型。}
  \label{tab:workflow_go_nogo}
  \begin{tabularx}{\textwidth}{@{}p{1.9cm}p{3.0cm}Y@{}}
    \toprule
    层级 & 继续前进前至少要确认的事 & 若未通过应如何处理 \\ \midrule
    PWEM & 候选模频段、简并标签和带图读取无歧义 & 回到单胞几何与归一化定义 \\
    CWT & 参数变化方向与主导耦合通道已识别 & 先缩小参数空间，不要急于全波 \\
    RCWA & 候选模开放谱与辐射排序稳定，功率守恒通过 & 回查傅里叶截断、材料色散和边界设定 \\
    FDTD/FEM & 阵列尺寸、PML、监视面和网格收敛，模式身份稳定 & 回查有限尺寸与开放边界，不要越级宣称器件优劣 \\
    电热光 & 注入、温升和工作点回写后目标模仍占优 & 不能输出最终器件窗口，应回到光学或注入分布层重设方案 \\
    \bottomrule
  \end{tabularx}
\end{table}

\section{2024--2026 前沿补充：代理模型与 AI 设计器如何接入互证链}
前沿工作显示，PCSEL 的快速设计工具正在从“低阶解析近似”扩展到“解析模型 + 数据驱动代理”的混合范式：一方面，快速工具评估研究给出不同近似器在精度/速度上的可操作对照\cite{lang2025rapidpcseltools}；另一方面，基于 Transformer 的代理模型已能在复杂晶格参数空间中实现高吞吐预测\cite{qi2025postsurrogate}。  
这并不改变本章的主结论，反而强化了“层级互证”的必要性：代理模型可以提前压缩搜索空间，但不能绕过物理保真层验收。

本书建议把代理模型接入点固定在“阶段 B 与阶段 C 之间”：
\begin{itemize}
  \item 先用代理模型做粗筛（候选区间、非可行区排除）；
  \item 再用 RCWA/FDTD/FEM 做物理验收（模式身份、边界条件、收敛与误差）；
  \item 最后用器件级电热模型验证工作窗口是否仍成立。
\end{itemize}
若缺少第二步和第三步，代理模型给出的“最优设计”只能算假设生成器，不应直接进入发布级结论。

\section{互证不是礼仪，而是发现错误的主要手段}
研究里常见的一种误区，是把多方法互证理解成“为了让论文更好看，多放几张图”。其实互证的主要价值是发现错误。原因很简单：不同方法对不同误差源敏感。若 RCWA 和 FDTD 对同一趋势意见一致，通常说明单胞开放物理和有限器件物理在该问题上相容；若它们结论冲突，就说明至少有一层物理被漏掉或被误设。

举例来说，若 CWT 预测某种椭圆孔扰动会增强偏振分裂，而 RCWA 也显示对应偏振通道的辐射损耗排序拉开，但 FEM 在真实有限外延层栈里却显示两支模式仍接近简并，那么问题很可能不在面内晶格，而在纵向层栈和边界。反之，若 RCWA 与 FEM 一致乐观，而器件级热电模型仍显示工作点下模式切换频繁，那么问题多半不在冷腔光学，而在注入与热反馈。

互证的作用不只是“确认”，还包括“定位”。

\begin{ExampleBox}
假设某方案在 RCWA 中显示目标模的光学阈值增益最低，于是看起来很理想；但 FEM 有限器件分析发现该模式在边缘附近场很强，导致边缘吸收一并增强；再进一步，热-电-光耦合分析又显示电流拥挤把增益集中在另一片区域。最终器件先起振的模式可能并不是 RCWA 中最优的那支。这不是 RCWA 错，而是它只负责单胞开放层级，后面两层信息把排序改写了。
\end{ExampleBox}

这个例子在于多方法链存在的意义：每一层都在回答自己的问题，任何一层都不应擅自越权宣布最终器件结论。

\section{什么时候该停在某一层，不再继续加方法}
并不是每个课题都需要把所有方法全部用满。若你的目标只是理解某种晶格扰动的基本耦合机制，CWT 加 RCWA 可能已经足够；若你的目标是论文级有限器件远场和偏振验证，通常还需要 FEM 或 FDTD；若你的目标是器件级阈值电流与高功率稳定工作区，那就必须继续进入后续的外延、电学和热学模型。

判断是否需要继续升级的一个实用标准是：当前层输出的信息，是否足以支持下一步设计决策。如果还不能，就必须升级；如果已经足够支持当前问题，不应为了“形式上更完整”而无条件加重模型。

\begin{PitfallBox}
“只要最后用了一次最重的全波仿真，前面的模型层级就都可以省略”是低效做法。这样会削弱结果的可解释性；一旦结果不理想，也很难判断应从哪个物理层级修改。对 PCSEL 这种跨层次器件，这类无分层路线尤其危险。
\end{PitfallBox}

\section{本章在后续器件章节前的作用}
本章的任务，是在进入 \cref{ch:epitaxial-optics,ch:qw-gain,ch:electrical,ch:eto-coupling} 之前，把光学方法层彻底整理清楚。因为后面的外延、增益、电注入和热效应章节都需要一个明确前提：哪些光学量已经可靠，哪些仍是近似代理量。如果这一点不先说清楚，后面的器件模型就会拿错输入，进而把整个链条的结论层级拉乱。

因此，本章既是数值方法部分的总结，也是器件物理部分的入口。它提醒读者：从这里开始，任何“器件结论”都必须建立在多层方法分工明确的基础上。

\section*{本章小结}
CWT、RCWA、FDTD 和 FEM 并不是彼此竞争的替代品，而是 PCSEL 建模链中分工不同的四个层次。CWT 负责机理压缩，RCWA 负责开放周期单胞精修，FDTD/FEM 负责有限器件全波验证。可靠结论来自这些方法的有序接力与互证，而不是某一种方法单独承担全部任务。对 PCSEL 这样同时包含周期光学、开放辐射和器件多物理场的对象，多方法链不是奢侈配置，而是基本要求。把这一点说得再直白一些：\cref{ch:bic} 中讨论的 BIC / quasi-BIC，并不是脱离数值方法的抽象修辞，它需要 RCWA 或 FEM 去分辨 $\Gamma$ 点附近辐射通道如何关闭或重新打开；而 \cref{ch:dynamics} 中讨论的调制、模式竞争和线宽，也不是脱离光学方法的后处理，它们必须继承这里已经建立的模式排序、损耗分解和器件输入接口。

\section*{练习题}
\begin{enumerate}
  \item \basicq 解释为什么“哪种方法更准”这个问题在没有说明问题层级时往往没有意义。

  \item \basicq 设计一条适用于大面积单模 PCSEL 的方法工作流，并说明每一步的淘汰标准。

  \item \midq 给出一个 RCWA 与 FDTD 结论可能冲突的场景，并说明冲突背后最可能缺失的物理是什么。

  \item \hardq 说明为什么互证时更应比较模式排序、对称性标签和参数趋势，而不是只比较单一绝对数值。
\end{enumerate}

\section*{建议继续阅读}
\begin{itemize}
  \item 下一章开始正式进入外延层的光学结构。后续章节将说明，前面光学方法的输出量如何被转换成外延设计中的纵向重叠、损耗重叠和上下辐射不对称。

  \item 若想补充方法基础，可回看 \cite{joannopoulos2008,taflove2005} 中关于本征法、散射法与时域法的相关章节。

  \item 若想检验自己是否理解了“方法层输出如何喂给后续模块”，可带着本章的分工框架回读 \cref{ch:bic,ch:dynamics}：前者依赖你正确分辨开放辐射通道，后者依赖你正确继承模式损耗和阈值排序。
\end{itemize}


